% Page 066 — page imprimée 63
% Suite de « Anecdotes, Témoignages, Souvenirs… »

Voilà donc ce récit conservé dans une famille, par la tradition orale. En fait il semble que si
cet assassinat a bien eu lieu, les pièces officielles, conservées à l’occasion du procès des
assassins les détails furent tout autres.

M. Maurice Abric de Fernouillet fut bien assassiné par trois hommes, le 16 janvier 1826 mais
au Château de l’Hom, une autre de ses propriétés, au Follaquier, hameau près de la route de
Saint André de Valborgne à quelques km des Salides et non pas aux Fons.
Sans doute y eut-il vol mais M. de Fenouillet, noble de fraîche date faisait preuve
d’arrogance, humiliant ses fermiers, corrigeant lui-même à coups de bâton deux enfants de 10
ans ! fils de son fermier, qui avaient dérobé quelques une de ses poires… n’hésitant pas à
couper les oreilles de chèvres qui avaient mangé ses châtaignes ou écorcé de jeunes
châtaigniers :

\begin{quote}
\textit{« …ce type là, raconte un habitant de Bassurel dans les années 1980, il aimait pas les gens.
Et alors les gens faisaient sortir leurs chèvres, comme ça. Et quand il trouvait des chèvres, il
leur coupait les oreilles ! Et quand ils ont vu ça, un jour ils l’ont tué, les voisins. Et je crois
qu’il est enterré là bas à l’Hom… Oh c’était une sale bête ! Il pouvait pas voir ses voisins,
quoi. »}
\end{quote}

Les trois assassins furent bel et bien condamnés à mort et guillotinés sur la place du
Pompidou le 7 mars 1827.

Cette triste histoire a tout de même causé la mort de quatre personnes, semé la désolation dans
plusieurs familles et laissé une dizaine d’orphelins dont certains n’ont pas connu leur père.

\begin{flushright}
Almanach de la Vallée Borgne 2006
\end{flushright}
